\chapter{Recomendaciones}\label{cap:recomendaciones}
Con base en la experiencia obtenida, se formulan las siguientes recomendaciones para trabajos futuros y para el uso educativo del proyecto.

\begin{enumerate}[leftmargin=2em,itemsep=4pt]
\item Incorporar entrenamiento por mini-lotes (\emph{mini-batch}) y optimizadores adaptativos como Momentum o Adam para acelerar la convergencia.
\item Añadir la función de pérdida de entropía cruzada con activación Softmax, más adecuada que el MSE para clasificación multiclase.
\item Ampliar el Experimento 3 a matrices de mayor resolución, a todo el alfabeto o a dígitos manuscritos, y separar un conjunto de validación.
\item Estudiar el efecto del número de neuronas ocultas y de la semilla aleatoria sobre la convergencia mediante experimentos repetidos.
\item Explorar la paralelización con los hilos de la biblioteca estándar de Rust para acelerar el entrenamiento de redes mayores.
\item Utilizar el repositorio como material didáctico en Programación 1, proponiendo a los estudiantes reproducir la prueba de escritorio y modificar los experimentos.
\item Para la versión final del documento, complementar la contextualización nacional y departamental con estadísticas oficiales verificables citadas en formato APA 7.
\end{enumerate}
